\chapter{Prospect Theory: Reference-Dependent Risky Choice}
\label{sec:p17-rationality-8-3}

Gl\"ockner's 2021 account distinguishes a descriptive theory of supplied
risky choices from a standard for how one ought to choose. Its organizing
questions are what prospect theory represents, how its versions differ,
where its predictions succeed or fail, and what choice data can establish
about cognitive processes \cite[pp.~485--493]{Glockner2021ProspectTheory}.
The subject is not exhausted by replacing the word ``utility'' in an
expected-utility formula. The reference point, the value transformation,
the weighting operation and the representation of the task all matter.

\section{Domain, standing, and four core features}

\begin{hbdefinition}{def:p17-pt-domain}{A represented risky choice}
The finite monetary domain here consists of lotteries
\(L=((o_1,p_1),\ldots,(o_m,p_m))\), with \(m\geq1\), monetary
outcomes \(o_i\in\mathbb R\), probabilities \(p_i\geq0\), and
\(\sum_i p_i=1\). Equal outcomes may be combined by adding their
probabilities; zero-probability branches may be removed. A reference
point \(r\in\mathbb R\) gives deviations \(x_i=o_i-r\): positive
deviations are gains, negative ones losses, and zero is neutral. A
value function \(v\) transforms deviations, while weighting functions
transform probabilities or cumulative probabilities. An evaluated
prospect receives a score \(V(L;r)\). Maximizing this score is a
conditional model prediction, including all maximizers when scores tie,
not a new normative obligation.

The usual description experiment provides the options, consequences and
numerical probabilities explicitly. An experience experiment provides
outcome samples instead. Neither format by itself supplies undiscovered
options, all relevant goals, or an account of real-world probability
learning \cite[pp.~485, 489--490]{Glockner2021ProspectTheory}.
\end{hbdefinition}

The source identifies four jointly important features. First, outcomes
are evaluated relative to a reference, not solely as final asset levels.
The reference can reflect an initial position, expectations or the way
the problem is represented. Second, diminishing sensitivity makes an
equal-sized increment less consequential farther from that reference:
the customary value curve is concave on gains and convex on losses.
Third, loss aversion gives a loss more negative value than an equally
large gain gives positive value, under the particular function and
magnitude comparison specified. Fourth, nonlinear weighting changes
the impact of probabilities; the characteristic inverse-S pattern
overweights small probabilities and underweights moderate or large ones.
The weighting function is not itself a subjective probability measure
\cite[pp.~485--486]{Glockner2021ProspectTheory}.

\begin{hbdefinition}{def:p17-pt-power-value}{A two-part power specification}
A standard parametric specification is
\[
 v(x)=
 \begin{cases}
  x^\alpha,&x\geq0,\\
  -\lambda(-x)^\beta,&x<0,
 \end{cases}
 \qquad \alpha,\beta>0,\quad\lambda>0.
\]
The customary diminishing-sensitivity case has
\(0<\alpha,\beta<1\). The gain and loss exponents describe different
curvatures; \(\lambda\) scales loss values. When \(\alpha=\beta\),
\(\lambda\) is the equal-magnitude loss/gain value ratio. Otherwise
that ratio is \(\lambda a^{\beta-\alpha}\) at magnitude \(a>0\),
and depends on the monetary scale and magnitude. This is a specified
function, not a claim that every person has the same parameters
\cite[p.~488]{Glockner2021ProspectTheory}
\cite[p.~309]{TverskyKahneman1992Cumulative}.
\end{hbdefinition}

\begin{hbconvention}{def:p17-pt-binary-weighting}{Binary evaluations used in the comparisons}
For the following binary comparisons, fix strictly increasing
\(w^+,w^-:[0,1]\to[0,1]\), with \(w^\pm(0)=0\) and
\(w^\pm(1)=1\). A prospect with gain \(b>0\) of probability
\(p\) and zero otherwise has score \(w^+(p)v(b)\); one with
loss \(-b\) of probability \(p\) and zero otherwise has score
\(w^-(p)v(-b)\). A sure deviation has its value \(v\).
The later cumulative definition extends this binary operation.
\end{hbconvention}

\begin{hbtheorem}{thm:p17-pt-binary-reference}{proved=annex:p17.rat8.3.binary-reference}{Curvature, weighting, and a separating reference change}
For the power specification with \(0<\alpha,\beta<1\), the value
function is increasing, concave for gains and convex for losses.
For a gain \(b>0\) occurring with probability \(p\in(0,1)\), and
zero otherwise, comparison with its sure monetary mean gives
\[
 \text{risky gain preferred to its sure mean}
 \quad\Longleftrightarrow\quad w^+(p)>p^\alpha.
\]
For a loss \(-b\) with probability \(p\), and zero otherwise,
\[
 \text{risky loss preferred to its sure mean}
 \quad\Longleftrightarrow\quad w^-(p)<p^\beta.
\]
Replace strict inequalities by equalities for indifference.

In the stipulated special case \(\alpha=\beta=1/2\),
\(\lambda=2\), and identity weighting, compare a certain final
outcome 4 with a lottery giving final outcomes 10 and 0, each with
probability one half. At reference 0 the certain outcome is preferred;
at reference 10 the lottery is preferred. A fixed expected-utility
function of the same final outcomes cannot reverse that comparison
merely because the reference label changes.
\end{hbtheorem}

These inequalities explain why the source's fourfold pattern needs
both transformations. Its characteristic pattern is risk aversion for
high-probability gains and risk seeking for high-probability losses,
with risk seeking for low-probability gains and risk aversion for
low-probability losses \cite[note~1, p.~491]{Glockner2021ProspectTheory}.
But \(w(p)>p\) alone does not imply \(w(p)>p^\alpha\) when
\(\alpha<1\). A proposed numerical model must pass the actual
comparison, rather than treating the pattern as a universal consequence
of the phrase ``small probabilities are overweighted.'' The reference
reversal above isolates reference-dependent value; its stipulated
parameters are not an empirical fit.

\section{Expected-value and expected-utility benchmarks}

The source's elementary choice is a certain payment of EUR~4 versus
EUR~10 with probability one half and zero otherwise. Expected money
is 4 versus 5. With the supplied power utility \(u(o)=o^{0.7}\) on
nonnegative final outcomes, however, the scores are approximately
2.639 and 2.506, so the certain option is preferred
\cite[p.~487]{Glockner2021ProspectTheory}. This is an example of
the different objectives represented by expected money and expected
utility, not a violation of expected utility.

\begin{hblemma}{thm:p17-pt-source-calculations}{proved=annex:p17.rat8.3.source-calculations}{The source's money examples and infinite-tail boundary}
For \(u(o)=o^a\), \(a>0\), the certain-4 option is preferred to
the half-chance-of-10 option exactly when
\(a<\log 2/\log(5/2)\), with a tie at equality.

In the source's St.~Petersburg convention, assume independent fair
coin flips. If \(K\) is the number of tails before the first head,
the payoff is zero for \(K=0\) and \(2^{K-1}\) for finite \(K\geq1\).
Set the payoff to zero on the probability-zero event of no head.
Its expected monetary payoff is infinite, whereas for
\(u(o)=o^a\), \(0<a<1\), its expected utility is
\[
                 \frac{1}{4(1-2^{a-1})}<\infty.
\]
This infinite-support illustration is separate from the finite CPT
domain above; no infinite-support CPT evaluation theorem is inferred.
\end{hblemma}

Gl\"ockner qualifies the expected-money benchmark by diminishing
marginal utility, additional goals, and nonrepeated decisions. Survival
or a minimum acceptable wealth level, joint outcomes, social comparison
and conformity may change what a utility function should represent.
Consequently a monetary expectation is not an adequate normative
criterion until the objective and domain are specified. Repetition
alone does not settle these questions
\cite[pp.~487--488, note~3]{Glockner2021ProspectTheory}.

The general expected-utility representation does \emph{not} require
concavity; the power function is the source's particular risk-averse
implementation. The exact conditional connection between concavity
and monetary risk aversion is
\cref{thm:p17-risk-aversion-concavity}. The finite-vNM theorem
\cref{thm:p17-finite-vnm} concerns complete, transitive, independent,
continuous preferences on its stated lottery domain, not the empirical
claim that every respondent obeys them. Nor does an arbitrary choice
of origin force the mathematical utility curve to be reflected at
that origin. The contrast with prospect theory is evaluation of
reference-coded deviations rather than a fixed utility of final
consequences, together with nonadditive decision weighting.

\section{Original and cumulative versions}

The 1979 original has an editing phase followed by evaluation. The
editing account codes gains and losses, combines probabilities of
identical outcomes, separates a certain component, removes shared
components in comparisons, simplifies small distinctions and detects
dominated alternatives. The original authors do not give a complete
context-independent ordering of all editing operations; one operation
can enable or prevent another. Its evaluation must therefore be read
with its editing assumptions intact
\cite[pp.~274--275]{KahnemanTversky1979Prospect}.

Write \(\pi\) for the original decision-weight function on
\([0,1]\), with \(\pi(0)=0\) and \(\pi(1)=1\); it need not
have \(\pi(p)+\pi(1-p)=1\).
For an edited regular prospect with at most two nonzero deviations,
\((x,p;y,q)\), remaining probability \(1-p-q\) at zero, the original
score is \(\pi(p)v(x)+\pi(q)v(y)\). ``Regular'' excludes a
strictly positive or strictly negative prospect of total nonzero
probability one. For the latter, with \(x\) more extreme than
\(y\), the original separate rule is
\[
             v(y)+\pi(p)\bigl(v(x)-v(y)\bigr).
\]
Here \(x>y>0\) or \(x<y<0\), and \(p+q=1\).
These are not interchangeable formulas. The original weighting scale
has formal properties, but not the later explicit parametric weighting
family; the source's reference to a missing formal description should
not be read as denying that the original paper defines decision
weights \cite[pp.~275--276]{KahnemanTversky1979Prospect}
\cite[p.~488]{Glockner2021ProspectTheory}.

\begin{hbdefinition}{def:p17-pt-cumulative}{Rank-dependent weights for supplied probabilities}
For this bounded risky-choice form of the 1992 cumulative model, use
strictly increasing functions \(w^+,w^-:[0,1]\to[0,1]\), each
with endpoints 0 and 1. Combine equal deviations, order the negative
ones from most negative to least negative, and the positive ones
from smallest to largest. A zero deviation contributes zero and is
omitted from the following sums. For negative deviation \(x_i\),
let \(P_i^-\) be the cumulative probability of deviations at least
as bad as \(x_i\), and \(P_{i-1}^-\) that of strictly worse ones.
For positive deviation \(x_i\), let \(P_i^+\) be the probability
of deviations at least as good as \(x_i\), and \(P_{i+1}^+\) that
of strictly better ones. Define
\[
 \pi_i^-=w^-(P_i^-)-w^-(P_{i-1}^-),\qquad
 \pi_i^+=w^+(P_i^+)-w^+(P_{i+1}^+),
 \qquad
 V(L;r)=\sum_{x_i<0}\pi_i^-v(x_i)+\sum_{x_i>0}\pi_i^+v(x_i).
\]
The outermost cumulative probabilities are zero. The reference and
functions are shared across the options in a comparison. This is a
probabilistic specialization, not the full 1992 representation for
uncertain events by capacities \cite[pp.~299--301]{TverskyKahneman1992Cumulative}.
\end{hbdefinition}

The chapter distinguishes five parameters in the customary 1992
specification: gain curvature \(\alpha\), loss curvature \(\beta\),
loss multiplier \(\lambda\), gain probability sensitivity
\(\gamma\), and loss probability sensitivity \(\delta\).
It expressly refers the full formalization to the 1992 paper
\cite[p.~488]{Glockner2021ProspectTheory}. The definition above
supplies the ranked finite-probability operation needed for the
chapter's own example without promising that separate paper's
axiomatic appendix or every uncertainty extension.

\begin{hbtheorem}{thm:p17-pt-coalescing}{proved=annex:p17.rat8.3.coalescing}{Weight mass, ranking, and splitting equal-outcome branches}
All defined cumulative weights are nonnegative. If \(P_+\) and
\(P_-\) are the total probabilities of strictly positive and
strictly negative deviations, their respective weight sums are
\(w^+(P_+)\) and \(w^-(P_-)\). A strictly one-sign prospect
of total mass one therefore has weights summing to one. For a
mixed prospect the combined nonzero-outcome weight need not be one.
Splitting one branch into several branches with the same outcome
and the same total probability leaves the score unchanged.

In the source's all-gain lottery with probabilities .20, .40, .40
and outcomes EUR~20, EUR~10, EUR~5, respectively, the weights are
\[
 w^+(.20),\qquad w^+(.60)-w^+(.20),\qquad 1-w^+(.60).
\]
Thus the two outcomes of probability .40 need not have equal
weights. For a stipulated increasing piecewise-linear function
through \((0,0),(.20,.30),(.60,.55),(1,1)\), they are .25 and
.45, while the top outcome has weight .30.
\end{hbtheorem}

The source's normalization language applies directly to its strictly
one-sign example; it must not be generalized to a mixed prospect.
Tversky and Kahneman explicitly permit mixed weight sums smaller
or greater than one \cite[p.~301]{TverskyKahneman1992Cumulative}.
Nor may \(w(p)\) be used as each individual outcome's weight in
CPT regardless of rank. Nonlinear weighting and cumulative weighting
are distinct operations \cite[note~4, p.~492]{Glockner2021ProspectTheory}.

The chapter also identifies weighting modifications reviewed by
Stott (2006), Fox--Tversky's (1998) two-stage account that first
estimates probabilities under uncertainty using support theory,
and Schmidt--Starmer--Sugden's (2008) lottery-valued reference points. These
extensions retain their separate sources and are not silently included
in the supplied-probability formula. The nearby ambiguity treatment
\cref{sec:p17-ambiguity-and-robust-choice} instead explicitly varies
the prior representation and decision criterion; unknown probabilities
are not resolved merely by distorting a known probability.

\section{Examples, predictive tests, and objections}

The chapter connects its four features to insurance and lottery
purchases, the equity premium, endowment effects, early sale of winning
stocks and retention of losing stocks, and late-day long-shot betting.
These are explanatory field illustrations attributed to its review,
not unique-mechanism identifications or prescriptions to buy, insure,
hold or sell \cite[p.~486]{Glockner2021ProspectTheory}.

Its framing illustration concerns equivalent descriptions of disease
programs. In the original classroom questionnaires, students at
Stanford and the University of British Columbia chose between saving
200 of 600 people for sure and a one-third chance of saving all 600.
For that gain frame, \(N=152\) and the reported percentages were
72\% for the sure program and 28\% for the risky one. A separate
loss-frame group, \(N=155\), chose between 400 deaths for sure and
a two-thirds chance of 600 deaths; 22\% chose the sure program and
78\% the risky one. The percentages are rounded source reports,
not population frequencies or an individual within-person reversal
\cite[p.~453]{TverskyKahneman1981Framing}.
The payoff distributions match because saved lives equal 600 minus
deaths. A representation-invariant preference over those consequences
must give the same comparison. Inferring irrationality still requires
that the two descriptions convey the same understood consequences,
probabilities and goals; the hypothetical classroom task does not
by itself validate medical or policy advice.

Gl\"ockner separates several empirical methods rather than treating
any successful fit as confirmation of the whole theory
\cite[pp.~488--491]{Glockner2021ProspectTheory}:

\begin{itemize}
\item \emph{Prediction outside the fitted observations.} His account of
Gl\"ockner--Pachur2012 compares incentivized decisions at two times,
with a diverse gamble set agreed in advance in an adversarial
collaboration. Parameters fitted at the first time predict the second.
The reported ordering favors CPT over EU, EV, the priority heuristic
and the other tested simple/adaptive heuristic models; fitting individual
parameters improves on the fixed 1992 parameters. This is evidence about
the tested simple, described, two-outcome lotteries and participants,
not all heuristics or all risk tasks. The chapter provides neither a
sample count nor confidence interval for this comparative summary;
none is inferred here. Selection of diagnostic problems in the 1979
paper and the risk of parameter overfitting motivate this comparison.
\item \emph{Critical properties that no parameter choice can remove.}
The chapter's Birnbaum illustration compares
\(A=(100,.85;50,.15)\) and \(B=(100,.95;7,.05)\).
It reports a majority preference for \(A\), but a reversal when
\(A\)'s lower branch is split into .10 and .05 at 50 and
\(B\)'s upper branch into .85 and .10 at 100. Each lottery's
distribution is unchanged. The coalescing theorem therefore says
that no refitting of the same distribution-based CPT functions can
account for that presentation reversal. Birnbaum's branch-sensitive
TAX alternative is a competing explanation, not proof of its universal
success \cite[pp.~463--465]{Birnbaum2008Paradoxes}.
\item \emph{More complex tasks.} Multi-outcome findings reviewed by
the chapter include choices better described by maximizing the total
probability of a winning outcome, and effects of time pressure and
environment structure on information search. They limit transfer
of simple-gamble predictive superiority. They do not prove that one
particular heuristic always dominates CPT in complex environments.
\end{itemize}

\section{Experience, sampling, and process evidence}

The chapter's experience example gives
\(A=(4,.8;0,.2)\) versus a certain 3. An illustrative sample is
\((4,0,4,4,4)\) from \(A\) and five 3s from the certain option.
The reported description--experience pattern favors the sure option
in description but the risky option in experience. The displayed
sample is an illustration, not a new participant record. In the
sampling paradigm exploration is followed by one paid final choice;
in the feedback paradigm each sampled outcome is paid. Those are
different incentive and learning problems
\cite[pp.~489--490, note~6]{Glockner2021ProspectTheory}.

\begin{hblemma}{thm:p17-pt-sampling}{proved=annex:p17.rat8.3.sampling}{The ten-sample calculation and a necessary qualification}
Under ten independent samples with rare-outcome probability .2,
the number \(N\) of rare outcomes is binomial with mean 2, and
\[
 \Pr(N<2)=.3758096384,\quad
 \Pr(N=2)=.3019898880,\quad
 \Pr(N>2)=.3222004736.
\]
Thus the chapter's rounded .38 and .32 values are correct for
this specified experiment. The inequality is not universal across
small sample sizes: with four independent samples at the same
probability, underrepresentation relative to mean .8 has probability
.4096, while overrepresentation has probability .5904.
\end{hblemma}

One debate asks whether experience requires qualitatively different
processes or can be explained by sampling-adjusted probability
estimates followed by prospect evaluation. The source reports an
ensemble winner in Erev et al.'s prediction competition, and subsequent
comparisons favorable to a two-stage CPT account in the sampling
paradigm. It also reports that the familiar direction of the gap is
not invariant: adjusted analyses and a broader task set can produce
\emph{more} small-probability overweighting in experience. Its
discussion of Gl\"ockner et al.2016, Kellen et al.2016 and
Wulff et al.2018 therefore retains problem dependence and competing
interpretations \cite[p.~490]{Glockner2021ProspectTheory}.
The ten-draw calculation is not a theorem that sampling alone
explains every observed gap, much less a prediction of every
experience participant's choices.

Finally, outcome prediction and process explanation are distinct.
The reviewed Mouselab studies require a mouse movement to reveal
information; complex tasks and time pressure change search patterns
and can induce adaptive strategy changes. Reviewed eye-tracking work
on simpler tasks finds broad information lookup, much search within
options, compensatory-looking choices, and fixations/response times
inconsistent with deliberate calculation of every transformed product.
These measurements have task-specific operational meanings: revealing
or fixating information is not identical to computing with it, and a
choice compatible with a formula does not establish conscious execution
of that formula. Evidence-accumulation and parallel-constraint models
are offered as candidate implementations, not identified uniquely by
these observations \cite[pp.~490--491]{Glockner2021ProspectTheory}.

\section{Use, comparison, and research boundary}

\begin{hbmethod}{def:p17-pt-use}{Derive and test a conditional risky-choice prediction}
Declare the available lotteries and what the person has actually
learned about them. Fix the reference-point convention, outcome
units, value specification, weighting version and parameter source
before predicting. Keep original editing assumptions separate from
the cumulative model. Combine equal outcomes for distribution-based
CPT, sort by sign and rank, compute cumulative weights and scores,
and report the complete maximizing set. For an empirical comparison,
separate fitting data from prediction data and state the task,
participant group, incentives, competing models, choice-noise
convention and uncertainty actually measured. If a needed reference,
parameter or learning model is unknown, report that identification
gap rather than fitting an explanation to the same response and
calling it a prospective prediction.

Use a separating case: the fixed-final-outcome reference reversal
tests what reference coding changes; the split-branch case tests
coalescing; the description/experience contrast tests the information
and learning conditions. A theory's success on one does not imply
success on the others.
\end{hbmethod}

The source regards predictive success as a reason to consider CPT
when reflecting on extended practical-rationality standards, while
acknowledging a different position in the heuristic literature. That
is an argumentative proposal requiring a bridge from predictive
adequacy and feasible goals to the proposed norm; the empirical
comparison alone does not supply that bridge
\cite[p.~489]{Glockner2021ProspectTheory}. The bounded-rationality
account in \cref{sec:p17-rationality-8-5} separately evaluates
resource limits and strategy/environment fit. It neither follows
from CPT's supplied-gamble formula nor disappears when CPT predicts
that particular task better.

Gl\"ockner's concluding 2021 assessment is strongest for simple
described risky choices, more qualified for complex tasks and
experience, and explicitly provisional about cognitive implementation.
Its research directions call for improved comparative tests,
experience/feedback-sensitive models and plausible integration
processes \cite[p.~491]{Glockner2021ProspectTheory}.
This assessment has descriptive and methodological standing. It
does not establish new rationality axioms, determine an individual's
parameters from one choice, or justify transferring a monetary-lottery
result to every consumer, social, clinical or policy decision.
